\section{Discussion}
\label{sec:discussion}
The main answer is conditional: within this parameterization, recipe and horizon, larger width makes delayed generalization more frequent and moves its held-out clock earlier, while leaving a long post-fit interval. Incidence, restricted onset, waiting and terminal success answer different questions. The learning-rate and multiplication controls extend the observed endpoint advantage, while low decay saturates the primary event clock and orbit splitting makes the conclusion threshold-dependent. A size recommendation therefore needs an explicit recipe, criterion and observation window.

P1 and P2 support the descriptive consequences of the candidate explanation. The much larger restricted-wait contrast than fit-clock contrast rules out an account solely in terms of earlier training fit. P3 nevertheless fails its planned fit-comparability gate. Frozen wide models can fit without generalizing, but their narrow counterparts cannot fit; hidden adaptation, shrinkage and optimization geometry remain confounded. The complete feature-learning explanation is unresolved, and Fourier or norm trajectories cannot select a unique mechanism. Further identifying interventions would need fit-comparable controls and a parameterization that separates those changes.

Budget choice also changes the empirical question. At a common update horizon, width can improve held-out learning; at 100 measured process-CPU seconds, wider runs have made fewer updates, and primary width~128 outperforms 512 descriptively. Low decay similarly has a large positive H-accuracy endpoint contrast but a small negative CPU-sampled one. These are distinct estimands, with unequal progress, early completion and one unavailable timing record explicitly retained. They establish neither an optimal width nor a hardware-independent scaling law.
